\SourceLeaf{411}%
\[
rr=cc+tt-2tc\cos B,\quad[rdr=(t-c\cos B)dt]
\]
\noindent [und]
\[
f=\beta+\frac{t(\gamma-\beta)}a;
\]
\noindent [folglich erhält man durch Integration der Formel für $d\log(r\sin\zeta)$ nach $t$ von $0$ bis $a$:]
\[
\begin{aligned}
\log\frac{b\sin C}{c\sin B}&=\tfrac1{24}\alpha(bb-cc)+\tfrac18\beta(bb-cc)\\
&\quad+\tfrac14(\gamma-\beta)\int\frac{tt-tc\cos B}a\,dt\\
&=\tfrac1{24}\alpha(bb-cc)+\tfrac18\beta(bb-cc)\\
&\quad+\tfrac1{24}(\gamma-\beta)(2aa-3ac\cos B)\\
&=\tfrac1{24}\alpha(bb-cc)+\tfrac18\beta(bb-cc)\\
&\quad+\tfrac1{48}(\gamma-\beta)(aa+3bb-3cc)\\
&=\tfrac1{24}\alpha(bb-cc)+\tfrac1{48}\beta(-aa+3bb-3cc)\\
&\quad+\tfrac1{48}\gamma(aa+3bb-3cc).
\end{aligned}
\]
\begin{flushright}W. Z. B. W.\end{flushright}
\par\Needspace{9\baselineskip}\begin{center}\textbf{Bemerkungen.}\end{center}
\noindent\hspace*{1.2em}Die Notiz [II] steht in einem Notizbuch kurz vor einer auf S. 444 abgedruckten Aufzeichnung mit dem Datum 1825. Dez. 4. Man darf daher annehmen, dass \textsc{Gauss} sie bei der Abfassung der unter [I] abgedruckten Briefe an \textsc{Olbers} und \textsc{Schumacher} niedergeschrieben hat. Den Inhalt des ersten Abschnitts der Notiz [II] hat \textsc{Gauss} im Jahr 1841 in \textsc{Crelle}s Journal veröffentlicht. Der Vollständigkeit halber wurde diese kleine, in den vierten Band der Werke nicht aufgenommene Abhandlung auch auf S. 451 bis 452 dieses Bandes abgedruckt.

In den \textit{Disquisitiones generales circa superficies curvas} § 28 sind die Winkelkorrekturen auf S. 402 wie folgt angegeben.
\begin{center}\begin{tabular}{lr}
Inselsberg & $\overset{\prime\prime}{4}{,}95131$\\
Hohehagen & $4{,}95113$\\
Brocken & $4{,}95104$\\
Summe & $14{,}85348$
\end{tabular}\end{center}
\noindent\hspace*{1.2em}Die Figur auf S. 403 wurde dem Text hinzugefügt.
\begin{flushright}Stäckel.\end{flushright}
